\clearpage\thispagestyle{empty}
\pdfbookmark[0]{Quellen zu Gauss an Bode (3. September 1820)}{bodeseptember1820-sources}
\begin{center}\textbf{Quellen zu Gauss an Bode (3. September 1820)}\end{center}
\noindent Diese KI-Übertragung ins heutige Deutsch folgt \textit{Gauss an Bode} in den \href{https://gdz.sub.uni-goettingen.de/id/PPN236007467}{\textit{Werken}, Band~VI} (Göttingen, 1874), S.~433--434. Der \href{https://api.digitale-sammlungen.de/iiif/image/v2/bsb10538352_00234/full/1600,/0/default.jpg}{Erstdruck} erschien unter \textit{Astronomische Beobachtungen im Jahr 1820} im von Johann Elert Bode herausgegebenen und 1820 in Berlin erschienenen \textit{Astronomischen Jahrbuch für das Jahr 1823}, S.~228--230. Die Links führen zum Exemplar der Bayerischen Staatsbibliothek. Der Brief beginnt unterhalb der vorhergehenden Tabelle. Gauß nennt von Staudt als Urheber der Berechnung der Pallas-Ephemeride, die im Brief erwähnt wird und außerhalb des in den Werken abgedruckten Textes steht.
\par\smallskip\noindent Für seine Beschreibung des Meridiankreises von Reichenbach verweist Gauß auf den \hyperlink{native-434-part-lower}{Bericht vom 5. Juni 1820}. Der \href{https://api.digitale-sammlungen.de/iiif/image/v2/bsb10538352_00236/full/1600,/0/default.jpg}{Erstdruck auf S.~230} schreibt den Namen Bradley. Die Werkausgabe schreibt Brandley. Diese Schreibweise ist in der Übertragung beibehalten. Die Tabellen folgen der Anordnung der Werkausgabe und behalten O und U für den oberen und unteren Sonnenrand bei. KI-Übertragung und KI-Satz.
\clearpage
